\chapter{LANDASAN TEORI}
\section{Tinjauan Pustaka}
\subsection{Konsep Dasar Aplikasi}

\vspace{0.5cm}
\noindent\fbox{\begin{minipage}{\textwidth}
\textit{\textbf{Pedoman Penulisan:} Definisi pengertian bahasa program yang digunakan.}
\end{minipage}}
\vspace{0.5cm}

\subsection{Metode Algoritma}

\vspace{0.5cm}
\noindent\fbox{\begin{minipage}{\textwidth}
\textit{\textbf{Pedoman Penulisan:} Definisi algoritma yang digunakan.}
\end{minipage}}
\vspace{0.5cm}

\subsection{Business Model Canvasing (BMC)}

\vspace{0.5cm}
\noindent\fbox{\begin{minipage}{\textwidth}
\textit{\textbf{Pedoman Penulisan:} Berisikan definisi BMC dan poin-poin yang terdapat dalam BMC.}
\end{minipage}}
\vspace{0.5cm}

\section{Penelitian Terkait}

\vspace{0.5cm}
\noindent\fbox{\begin{minipage}{\textwidth}
\textit{\textbf{Pedoman Penulisan:} Analisa tentang aplikasi yang akan dibahas dibanding penelitian terdahulu.}
\end{minipage}}
\vspace{0.5cm}

\section{Teori Pendukung}
\subsection{Pengujian Aplikasi}

\vspace{0.5cm}
\noindent\fbox{\begin{minipage}{\textwidth}
\textit{\textbf{Pedoman Penulisan:} Contoh: Teori Blackbox testing.}
\end{minipage}}
\vspace{0.5cm}

\subsection{Peralatan Pendukung}

\vspace{0.5cm}
\noindent\fbox{\begin{minipage}{\textwidth}
\textit{\textbf{Pedoman Penulisan:} Contoh: Teori Android Studio, Java, dll.}
\end{minipage}}
\vspace{0.5cm}

